% !TEX root = ../HDR.tex
%%%%%%%%%%%%%%%%%%%%%%%%%%%%%%%%%%%%%%%%%%%%%%%%%%%%%%%%%%%%%%%%%%%%%%%%%%%%%%%%%%
\chapter{Construction de connaissance artificielle}

\section{L'autonomie artificielle}

Définir l'autonomie des agents artificiels n'est pas aussi facile qu'il n'y parait. 
David Vernon \cite{vernon_artificial_2014} identifie plusieurs types d'autonomie

\section{Théorie de l'enaction}

Francisco Varella et humberto Maturana \cite{varela_embodied_1993} sont à l'origine de la théorie de l'enaction.

Tom froese et Tom Ziemke ont propose l'IA enactive 
\cite{froese_enactive_2009}


\section{Théoride des contingences sensorimotrices}

Kevin O'Regan et Alva Noë ont proposé la théorie des contingences sensorimotrices \cite{oregan_sensorimotor_2001}

\section{L'émergence de sémantique ancrée dans l'expérience sensorimotrice}

Pei Wang \cite{wang_experience-grounded_2005} pose la question de l'émergence sémantique ancrée dans l'expérience sensorimotrice 


\section{L'apprentissage des structures spatiales}

\cite{gay_autonomous_2017,terekhov_learning_2019,raju_space_2022}

\section{Les architectures cognitives bio-inspirées}

Ron Sun \cite{sun_desiderata_2004} pose les bases des attentes envers les architectures cognitives.

David Vernon \cite{vernon_desiderata_2016} rebondit sur l'article fondateur de Sun pour spécifier les désiderata pour une architecture cognitive développementale. 

\subsection{Constructivist Anticipatory Learning Mechanism (CALM)}

\cite{perotto_computational_2013}

\subsection{LIDA}

Sean Kugele \cite{kugele_constructivist_2025} étend LIDA pour inclure un apprentissage constructiviste. 

\subsection{AERA}

L'équipe de Reykjavik développe AERA 